% discussion.tex -- Sequencer Paper (BNR-2026-009)

\section{Discussion}

\subsection{Key Findings}

\textbf{Finding 1: Single-operator sequencing provides $>$999~ms headroom per block.}
At the maximum tested load of 5{,}000 transactions per block, the sequencer's pure
overhead (mempool drain + block assembly) consumes $<$1~ms of the 1{,}000~ms block
interval. This leaves $>$999~ms for EVM execution. Based on published Geth benchmarks
of 3{,}000--5{,}000 simple transfers per second~\cite{geth2024tracing}, a 1-second
block can accommodate 3{,}000--5{,}000 transactions with EVM execution, placing the
sequencer overhead at $<$0.1\% of the total block time budget. The sequencer is not
and will not be the bottleneck.

\textbf{Finding 2: FIFO ordering is deterministic under all tested conditions.}
Across 15 benchmark configurations, including concurrent access with 4 producer
goroutines, adversarial censoring scenarios, and high-throughput stress tests with
250{,}000 pre-loaded transactions, FIFO ordering accuracy remains at exactly 100\%.
The sequence number mechanism (atomic counter at insertion time) provides a total
order that is preserved through the mempool drain and block assembly pipeline.
This determinism simplifies downstream ZK proving: the prover does not need to
verify ordering correctness, only state transition validity.

\textbf{Finding 3: Forced inclusion achieves single-tick latency.}
Under both cooperative and adversarial conditions, all forced transactions are included
within one block tick (${\sim}$1~second) of their deadline expiry. The maximum observed
latency is 1{,}010.6~ms. This is three orders of magnitude faster than the 24-hour
deadline window, providing substantial safety margin. In practice, the forced inclusion
mechanism drains the queue at every block production cycle, meaning forced transactions
are never deferred beyond the first block after their deadline.

\textbf{Finding 4: Concurrent access degrades insertion throughput but preserves correctness.}
Single-threaded mempool insertion achieves ${\sim}$2.8M~tx/s, while concurrent access
with 4 goroutines reduces this to ${\sim}$33K~tx/s due to mutex contention. However,
33K~tx/s still exceeds the target throughput (500~TPS) by a factor of 66$\times$, and
FIFO ordering remains perfect. For production use, the mutex-based approach is sufficient;
lock-free alternatives (e.g., MPSC channels) are available as optimization options if
contention becomes measurable.

\textbf{Finding 5: The enterprise context eliminates MEV as a design concern.}
The zero-fee transaction model, absence of public DEX contracts, and single known
enterprise operator collectively eliminate the economic incentives for MEV extraction.
FIFO ordering provides fairness without the complexity of MEV-mitigation mechanisms
(threshold encryption, commit-reveal schemes, fair ordering protocols) required in
public L2 systems.

\subsection{MEV Analysis}

We analyze the three primary MEV attack vectors and their applicability to the Basis
Network enterprise context:

\begin{table}[htbp]
\caption{MEV Attack Vector Analysis}
\label{tab:mev-analysis}
\centering
\begin{tabular}{lll}
\toprule
\textbf{Attack} & \textbf{Prerequisite} & \textbf{Basis Status} \\
\midrule
Front-running & Fee priority, public mempool & Eliminated (FIFO) \\
Sandwich & DEX with slippage & No public DEXs \\
Liquidation & DeFi lending protocols & No DeFi protocols \\
Time-bandit & Block reorg incentive & Single operator \\
Censoring & Sequencer discretion & Forced inclusion \\
\bottomrule
\end{tabular}
\end{table}

The only MEV vector that applies in the enterprise context is transaction censoring,
which is addressed by the forced inclusion mechanism. The remaining vectors require
infrastructure (DEXs, lending protocols, fee markets) that does not exist on
enterprise-operated chains.

\subsection{Invariants Discovered}

The following invariants are established by this investigation and must be preserved
by downstream pipeline stages (Logicist, Architect, Prover):

\begin{enumerate}
  \item \textbf{INV-SEQ-1 (FIFO Ordering)}: The sequencer SHALL order transactions
    strictly by arrival time within each block. No priority ordering, no fee-based
    reordering. Transactions from the forced inclusion queue are prepended before
    mempool transactions but maintain FIFO order within each source.

  \item \textbf{INV-SEQ-2 (Forced Inclusion Liveness)}: Every transaction submitted
    to the forced inclusion queue on L1 SHALL be included in an L2 block within
    the configured deadline. The sequencer MUST drain the forced inclusion queue at
    the start of every block production cycle, before processing mempool transactions.

  \item \textbf{INV-SEQ-3 (Block Production Regularity)}: The sequencer SHALL produce
    blocks at fixed intervals ($\Delta t = 1$~second). Block production MUST NOT be
    skipped when the mempool is empty---empty blocks maintain the chain's liveness
    guarantee and simplify downstream batch aggregation.

  \item \textbf{INV-SEQ-4 (Mempool Capacity Bound)}: The mempool SHALL enforce a
    maximum capacity $C_{max}$. When full, new insertions are rejected with an
    explicit error (not silently dropped). This bounds memory consumption at
    $C_{max} \times 168$~bytes for the queue structure.

  \item \textbf{INV-SEQ-5 (Single-Operator Trust Model)}: The sequencer operates under
    a trusted-but-verified model. The operator is trusted to include transactions
    fairly (FIFO), but the forced inclusion mechanism provides a cryptographic
    fallback: if the operator censors a transaction, the user can bypass the sequencer
    entirely via L1 submission. The ZK proof verifies the resulting state transition,
    not the ordering policy.
\end{enumerate}

\subsection{Forced Inclusion Contract Design}

The L1 forced inclusion contract follows the Arbitrum DelayedInbox pattern:

\begin{lstlisting}[language=Go, caption={Forced Inclusion L1 Contract Interface (Solidity-like pseudocode)}, label={lst:forced-contract}]
// User submits tx to L1 queue
function submitForcedTx(
  bytes calldata txData
) external {
  queue.push(ForcedTx({
    data: txData,
    sender: msg.sender,
    deadline: block.timestamp + DEADLINE,
    index: queue.length
  }));
}

// Anyone can force after deadline
function forceInclusion(
  uint256 index
) external {
  require(
    block.timestamp >= queue[index].deadline
  );
  // emit event for sequencer to process
}
\end{lstlisting}

The contract maintains a FIFO queue indexed by submission order. The sequencer monitors
L1 events and includes pending forced transactions before their deadline. After the
deadline, any user can trigger \texttt{forceInclusion()}, which emits an event that
the sequencer must process in the next block.

\subsection{Residual Risks}

\textbf{EVM execution time.} Block production headroom ($>$999~ms) must accommodate
EVM execution. At 500~TPS with 21K~gas/tx simple transfers, Geth processes
${\sim}$10{,}000~TPS---providing a 20$\times$ safety margin. Complex contract
interactions (higher gas per transaction) reduce this margin but remain well within
budget for enterprise workloads.

\textbf{L1 monitoring latency.} The forced inclusion queue polls L1 for events.
Avalanche's sub-second finality ensures that L1 events are visible within 1--2 seconds,
but network latency and RPC polling intervals add variable delay. This latency is
bounded by the 24-hour deadline and does not affect correctness.

\textbf{State growth.} The mempool uses 168~bytes per transaction in the queue
structure. At 500~TPS sustained, mempool memory grows at ${\sim}$84~KB/s. With a
capacity limit of 100{,}000 transactions, peak memory is ${\sim}$16~MB---negligible
for server hardware.

\textbf{Windows timer resolution.} Some direct measurements report 0~$\mu$s due to
Windows timer granularity (${\sim}$15~ms). The Go native benchmarks resolve this by
using the \texttt{testing.B} framework with nanosecond precision. Production benchmarks
should be conducted on Linux for consistent timer behavior.

\subsection{Limitations}

\textbf{No EVM execution.} The prototype measures pure sequencer overhead (mempool
drain + block assembly) without EVM transaction execution. This is by design: EVM
execution is the domain of RU-L1 (BNR-2026-008). The combined system performance
requires integration testing.

\textbf{No network layer.} Transaction arrival in the prototype is via direct function
call, not RPC or P2P networking. Network latency and serialization overhead will add
to end-to-end latency but do not affect sequencer-internal performance.

\textbf{No persistent storage.} Blocks are held in memory, not written to disk. WAL
(Write-Ahead Log) and database persistence will add per-block overhead. The validium
node's WAL (benchmarked in RU-V4) adds ${\sim}$0.1~ms per write, which is negligible.

\textbf{Single-machine benchmarks.} All measurements are on a single Windows 11
machine with 20 cores. Production deployment on Linux server hardware may show
different absolute numbers but the relative relationships and scaling behavior
are expected to hold.
